\section{Open problems}\label[appendix]{sec:A}
\textbf{The one-sorted coding arity, settled.} In the first version of this work the value of ca(SE₂) was left open at $n = 2$. It is now determined: \cref{corollary:4.4} gives ca(SE₂) = 1. No single quadratic equation with any finite number of integer unknowns defines the graph of an injection $ℤ^{2} → ℤ$. The mixed-semantics analogue, with inputs and output in ℕ and unknowns in ℤ, is a different question with a different answer --- it equals 2, witnessed by the Cantor pairing (\cref{sec:B}) --- and the two must not be identified. That they genuinely differ is itself a theorem (\texttt{natInput\_intInput\_separation\_at\_two}): \texttt{AtomGraphZ 2 2} and $¬$\,\texttt{AtomGraphZZ 2 2} hold simultaneously. In particular \cref{theorem:4.1} itself does not survive replacing ℤ by ℕ, and neither does the symmetric-point collision on a full lattice coset that proves it; it is the theorem and that collision that fail, not every individual step.

\textbf{A one-witness normal form.} For a single quadratic equation with exactly one unknown one expects a finite-union normal form of its fibres, a different object from the normal forms of total functions in \cref{sec:10}: a finite exceptional set, finitely many polynomial progressions of degree at most two, and finitely many affine projections of unit orbits with the exponent restricted to a congruence class. Establishing it requires the finiteness of the number of unit-orbit classes for x² − Dy² = N in orbit-decomposition form. We state the expectation and do not prove it here; a draft formalization is in the development notes accompanying the artifact, where it is marked as conditional and unproved, and nothing in Sections 1–9 uses it.

\textbf{Effectivity.} The constant K of \cref{theorem:9.2}, and hence the constants of \cref{theorem:2.1,theorem:2.2,theorem:3.2}, are not effective. \cref{theorem:4.1} carries no constant: its coset branch is a collision argument, not a count. Everything else in those proofs is explicit, so an effective version of \cref{theorem:9.2} would make all of them effective with constants that can be read off mechanically. We do not claim that the small-solution bounds cited in \cref{sec:11} yield such a version: those bound solutions of a single equation, whereas what is needed is a bound uniform over all affine perturbations of a fixed quadratic part.

\textbf{Beyond \cref{sec:10}.} Three questions are not addressed here. A Lean proof of $\mathrm{TCER}$ itself would remove the one analytic hypothesis of \cref{sec:10}; the artifact states it as a hypothesis and \cref{lemma:10.8} derives it from published work. The least number of unknowns with which a given SE₂-definable function can be defined is not determined in general; \cref{proposition:10.18} settles one example. And the classification of \cref{sec:10} concerns graphs of total functions: a single quadratic equation defines more general relations, including the relation $a ≥ 0$ through four squares, and those are not classified.
